
This paper introduced the YouRA evaluation framework and reported pipeline validation results for the hypothesis that calibration quality, hallucination resistance, and adversarial robustness in LLMs co-vary along a common latent dimension — \textit{epistemic reliability} — that is largely independent of raw MMLU capability.

Under a synthetic score matrix conforming to the hypothesized latent structure (N=30 models, 5 trustworthiness metrics):

\begin{enumerate}
\item \textbf{The hypothesized latent structure is recoverable.} Partial ρ(ECE, TruthfulQA\%|MMLU) = −0.758 (BCa 95\% CI [−0.894, −0.504]); all 10 cross-metric partial correlations exceed |ρ| = 0.40; Factor 1 explains 72.1\% of shared variance (KMO = 0.879, Tucker's congruence = 1.000).
\end{enumerate}

\begin{enumerate}
\item \textbf{The calibration–hallucination correlation is not attributable to MMLU capability.} Survival fraction = 0.943; MMLU control reduces the correlation by 5.7\%. Discriminant validity with HumanEval: |partial ρ| = 0.082.
\end{enumerate}

\begin{enumerate}
\item \textbf{Incremental predictive power over capability-only screening is uncertain.} LOO-AUC = 0.739; ΔAUC = 0.051 (BCa 95\% CI [−0.194, 0.449]). The CI width indicates insufficient power at N=30 to resolve whether a meaningful incremental advantage exists.
\end{enumerate}

\begin{enumerate}
\item \textbf{The mechanistic pathway via embedding perturbation (H-M3) was not executed.} This remains an open empirical question.
\end{enumerate}

All quantitative results are synthetic-data pipeline validation. Real-data replication via lm-evaluation-harness on the target 30-model population is the immediate scientific priority before any claim about actual LLM behavior can be advanced.

